\documentclass[UTF8,10pt,letterpaper,fontset=fandol]{ctexart}
\usepackage[margin=1in]{geometry}
\usepackage{amsmath,amssymb,enumitem,needspace}
\usepackage[unicode,hidelinks]{hyperref}
\pagestyle{plain}
\raggedbottom
\setlength{\parindent}{0pt}
\setlength{\parskip}{1pt}
\setlength{\emergencystretch}{2em}
\linespread{1.1}
\AtBeginDocument{%
\setlength{\abovedisplayskip}{6pt}%
\setlength{\belowdisplayskip}{6pt}%
\setlength{\abovedisplayshortskip}{3pt}%
\setlength{\belowdisplayshortskip}{3pt}}
\setcounter{secnumdepth}{0}
\setlist[enumerate]{label=(\arabic*),leftmargin=2em,itemsep=0.25em,topsep=0.2em}
\ctexset{section={format=\Large\bfseries,beforeskip=1.2ex,afterskip=0.8ex}}
\newcommand{\C}{\mathbb C}
\newcommand{\F}{\mathbb F}
\begin{document}
\begin{center}
{\LARGE Homework 7}\\[0.6em]
October 28, 2026\\[0.25em]
Name: YuanHao Tang\quad Student ID: 3120265621
\end{center}

\section{Problem 1}
Let $\mathcal V$ be a linear space over $\F$, $\mathcal A\in L(\mathcal V)$, $\mathcal W\subset\mathcal V$ be $\mathcal A$-invariant, and $\psi(\lambda)$ be the minimal polynomial of $\mathcal A$ at $\mathcal W$. Show that there exists $x\in\mathcal W$ such that the minimal polynomial of $\mathcal A$ at $x$ is $\psi(\lambda)$.

\par\Needspace{4\baselineskip}\textbf{Solution}\par
If $\mathcal W=\{0\}$, take $x=0$ and $\psi=1$. Otherwise write $\psi=\prod_{i=1}^r p_i^{m_i}$ with distinct monic irreducibles $p_i$, and put $q_i=\psi/p_i^{m_i}$. Minimality of $\psi$ provides $v_i\in\mathcal W$ with $(\psi/p_i)(\mathcal A)v_i\ne0$. Set $y_i=q_i(\mathcal A)v_i\in\mathcal W$. Then

\[
p_i(\mathcal A)^{m_i}y_i=0,\qquad p_i(\mathcal A)^{m_i-1}y_i\ne0.
\]
Thus the minimal polynomial at $y_i$ is $p_i^{m_i}$. These polynomials are pairwise coprime, so the coprime-sum property gives

\[
\boxed{x=\sum_{i=1}^r y_i\in\mathcal W,\qquad m_x=\prod_{i=1}^r p_i^{m_i}=\psi}.
\]

\section{Problem 2}
Let $\mathcal A\in L(\C_n[\lambda])$ be defined as

\[
\mathcal A[f(\lambda)]=f(\lambda+a),\qquad\forall f(\lambda)\in\C_n[\lambda].
\]
for a nonzero constant $a\in\C$; its matrix representation $A$ with respect to the basis $1,\lambda,\lambda^2,\ldots,\lambda^n$ was found in Homework 5. Find out

\begin{enumerate}
\item the Smith canonical form $S$ of the eigenmatrix $\lambda E-A$, and the Jordan canonical form $J$ of $A$;

\item the transition matrix $X$ that reduces $A$ to the Jordan canonical form $J$, i.e. $X^{-1}AX=J$.

\end{enumerate}

\par\Needspace{4\baselineskip}\textbf{Solution}\par

\begin{enumerate}
\item Write the polynomial variable as $s$ and put $\Delta=\mathcal A-I$. It lowers nonconstant polynomial degrees, with $\Delta^{n+1}=0$ and $\Delta^n(s^n)=n!a^n\ne0$. Hence there is one Jordan block of size $n+1$:

\[
\boxed{S=\operatorname{diag}(I_n,(\lambda-1)^{n+1}),\qquad J=J_{n+1}(1)}.
\]

\item Take the following polynomials and use their coordinates in $(1,s,\ldots,s^n)$ as the columns of $X$, indexed from zero:

\[
f_0(s)=1,\qquad
f_k(s)=\frac{1}{k!a^k}\prod_{r=0}^{k-1}(s-ra)\quad(1\le k\le n).
\]

\[
\boxed{X_{ik}=[s^i]f_k(s)\quad(0\le i,k\le n)},\qquad
X_{kk}=\frac{1}{k!a^k}\ne0.
\]
Here $[s^i]$ extracts a coefficient. The upper triangular $X$ is invertible. Since $f_k(s+a)=f_k(s)+f_{k-1}(s)$, with $f_{-1}=0$, we have $AX=XJ$ and hence $X^{-1}AX=J$.

\end{enumerate}

\newpage

\section{Problem 3}
Let $E$ be an identity matrix with appropriate dimension,

\[
\begin{gathered}
\lambda+3,\quad\lambda+3,\quad(\lambda+3)^2,\quad\lambda+2+5j,\quad\lambda+2-5j,\\
(\lambda+2+5j)^2,\quad(\lambda+2-5j)^2.
\end{gathered}
\]
are elementary divisors for $\lambda E-A$.

\begin{enumerate}
\item Please find out the Smith canonical form $S$ for $\lambda E-A$, the natural normal form $N$ and the Jordan canonical form $J$ for $A$.

\item Let $\lambda_1=-3$, $\lambda_2=-2+5j$, $\lambda_3=-2-5j$ and

\[
T(\lambda_i)\triangleq\{x:\exists k\in\mathbb Z_+,\ (A-\lambda_iE)^kx=0\}.
\]
$\mathbb Z_{+,0}$ is the set of non-negative integers. Show that $T(\lambda_i)=\ker[(A-\lambda_iE)^2]$, $i=1,2,3$, and determine the corresponding dimensions $\dim T(\lambda_i)$, $i=1,2,3$.

\end{enumerate}

\par\Needspace{4\baselineskip}\textbf{Solution}\par
Use $j^2=-1$. Let $J_m(\rho)$ have diagonal $\rho$ and superdiagonal ones; $C(p)$ is the column companion matrix, with subdiagonal ones and negated coefficients of monic $p$ in ascending order in its last column.

\begin{enumerate}
\item Put $x=\lambda+3$, $q=(\lambda+2)^2+25$. Sorting exponents and padding on the left gives nonconstant invariants $x,xq,x^2q^2$, of degrees one, three and six. Thus

\[
\boxed{S=\operatorname{diag}(I_7,x,xq,x^2q^2),\qquad N=\operatorname{diag}(-3,C(xq),C(x^2q^2))}.
\]
Write $K_\pm(a,b)=\operatorname{diag}(a\pm bj,J_2(a\pm bj))$. Then

\[
\boxed{J=\operatorname{diag}(-3,-3,J_2(-3),K_-(-2,5),K_+(-2,5))}.
\]

\item At each eigenvalue the largest Jordan block has size two, so $(A-\lambda_iI)^2$ vanishes on its generalized eigenspace. The operator $A-\lambda_iI$ is invertible on all other components. Therefore

\[
\boxed{T(\lambda_i)=\ker(A-\lambda_iI)^2,\qquad (\dim T(\lambda_1),\dim T(\lambda_2),\dim T(\lambda_3))=(4,3,3)}.
\]

\end{enumerate}

\newpage

\section{Problem 4}
Given a matrix $A$ whose elementary divisor are

\[
\begin{gathered}
(\lambda+1),\quad(\lambda+1)^2,\quad(\lambda+1+j),\quad(\lambda+1-j),\\
(\lambda+1)^2,\quad(\lambda+1+j)^2,\quad(\lambda+1-j)^2.
\end{gathered}
\]

\begin{enumerate}
\item Find out the Smith canonical form $S$ for $\lambda E-A$;

\item Find out the Natural normal form $N$, the complex Jordan canonical form $J_{\rm complex}$ and the real Jordan canonical form $J_{\rm real}$ for $A$;

\item The eigenvalues of matrix $A$ are denoted as $\lambda_1=-1$, $\lambda_2=-1-j$, $\lambda_3=-1+j$. For $\lambda_i$, $i=1,2,3$, give the geometric multiplicity $q_i$ and algebraic multiplicity $\widehat n_i$;

\item Find out $\dim T(\lambda_i)$, $i=1,2,3$, where

\[
T(\lambda_i)\triangleq\{x:\exists k\in\mathbb Z_+,\ (A-\lambda_iE_n)^kx=0\}.
\]

\end{enumerate}

\par\Needspace{4\baselineskip}\textbf{Solution}\par
Put $x=\lambda+1$, $q=x^2+1$ and retain the block notation from Problem 3.

\begin{enumerate}
\item The nonconstant invariants are $x,x^2q,x^2q^2$, of total degree eleven, so

\[
\boxed{S=\operatorname{diag}(I_8,x,x^2q,x^2q^2)}.
\]

\item
\[
\boxed{N=\operatorname{diag}(-1,C(x^2q),C(x^2q^2))},
\]

\[
\boxed{J_{\rm complex}=\operatorname{diag}(-1,J_2(-1),J_2(-1),K_-(-1,1),K_+(-1,1))}.
\]
Put $R=\begin{pmatrix}-1&1\\-1&-1\end{pmatrix}$ and $R_2=\begin{pmatrix}R&I_2\\0&R\end{pmatrix}$. Pairing conjugate blocks gives

\[
\boxed{J_{\rm real}=\operatorname{diag}(-1,J_2(-1),J_2(-1),R,R_2)}.
\]

\item Geometric multiplicity counts blocks, while algebraic multiplicity sums their sizes:

\[
\begin{array}{c|cc}\lambda_i&q_i&\widehat n_i\\\hline -1&3&5\\-1-j&2&3\\-1+j&2&3\end{array}
\]

\item Generalized eigenspace dimensions equal algebraic multiplicities, hence

\[
\boxed{\dim T(\lambda_1)=5,\qquad\dim T(\lambda_2)=3,\qquad\dim T(\lambda_3)=3}.
\]

\end{enumerate}

\end{document}
